\section{Билет 34. Основные положения молекулярно-кинетической теории. Характерные массы вещества. Количество вещества. Число Авогадро. Число Лошмидта.}

\begin{definition}
\textbf{Молекулярно-кинетическая теория (МКТ)} --- раздел физики, изучающий строение и свойства веществ на основе представлений о том, что все тела состоят из мельчайших частиц (молекул и атомов), находящихся в непрерывном хаотическом движении.
\par\noindent\textbf{Три основных положения МКТ:}
\begin{enumerate}
	\item \textbf{Дискретность. }Все вещества --- жидкие, твёрдые и газообразные --- образованы из мельчайших частиц: молекул, которые сами состоят из атомов. Молекулы и атомы электрически нейтральны.
	\item \textbf{Тепловое движение. }Атомы и молекулы находятся в непрерывном хаотическом тепловом движении.
	\item \textbf{Взаимодействие. }Частицы взаимодействуют друг с другом силами электрической природы; гравитационное взаимодействие между ними пренебрежимо мало.
\end{enumerate}
\end{definition}

\begin{definition}
Для характеристики масс атомов и молекул используются следующие величины:
\begin{itemize}
    \item \textbf{Относительная атомная масса} элемента $A$ --- отношение массы атома этого элемента к $1/12$ массы атома изотопа углерода $^{12}\text{C}$:
    \begin{equation}
        A = \frac{m_X}{\frac{1}{12} m_{^{12}\text{C}}}.
    \end{equation}
    Значения $A$ приведены в Периодической системе Д.И. Менделеева.
    \item \textbf{Относительная молекулярная масса} вещества $M_r$ --- отношение массы молекулы к $1/12$ массы атома $^{12}\text{C}$:
    \begin{equation}
        M_r = \sum_i n_i A_i,
    \end{equation}
    где $n_i$ --- число атомов $i$-го элемента в молекуле, $A_i$ --- его относительная атомная масса.
    \item \textbf{Атомная единица массы (а.е.м.)} --- $1\,\text{а.е.м.} = \frac{1}{12} m_{^{12}\text{C}} \approx 1.66 \cdot 10^{-27}$ кг.
\end{itemize}
\end{definition}

\begin{definition}
\textbf{Количество вещества} $\nu$ --- физическая величина, определяемая числом структурных элементов (атомов, молекул, ионов) в системе. Единица измерения --- моль.
\begin{equation}
    \nu = \frac{N}{N_A},
\end{equation}
где $N$ --- число частиц, $N_A$ --- \textbf{постоянная Авогадро}:
\begin{equation}
    N_A = 6.02 \cdot 10^{23} \text{ моль}^{-1}.
\end{equation}
Один моль любого вещества содержит $N_A$ структурных элементов.
\end{definition}

\begin{definition}
\textbf{Молярная масса} $M$ --- масса одного моля вещества:
\begin{equation}
    M = N_A \cdot m_0,
\end{equation}
где $m_0$ --- масса одной молекулы. Молярная масса выражается в кг/моль. Связь с относительной молекулярной массой:
\begin{equation}
    M = M_r \cdot 10^{-3} \text{ кг/моль}.
\end{equation}
Количество вещества через массу $m$ и молярную массу:
\begin{equation}
    \nu = \frac{m}{M}.
\end{equation}
\end{definition}

\begin{property}[Число Лошмидта]
\textbf{Число Лошмидта} $N_L$ --- число молекул, содержащихся в единице объёма идеального газа при нормальных условиях ($p_0 = 1.01 \cdot 10^5$ Па, $T_0 = 273$ К):
\begin{equation}
    N_L = \frac{p_0}{k T_0} = 2.68 \cdot 10^{25} \text{ м}^{-3},
\end{equation}
где $k = 1.38 \cdot 10^{-23}$ Дж/К --- постоянная Больцмана.
\par\noindent\textbf{Следствие:} При одинаковых температуре и давлении все газы содержат в единице объёма одинаковое число молекул (закон Авогадро).
\end{property}

\begin{remark}
\textbf{Молярный объём} $V_M$ --- объём одного моля газа при нормальных условиях:
\begin{equation}
    V_M = \frac{V}{\nu} = 22.4 \cdot 10^{-3} \text{ м}^3/\text{моль} = 22.4 \text{ л/моль}.
\end{equation}
Это значение одинаково для всех идеальных газов, что подтверждает справедливость уравнения состояния $pV = \nu RT$.
\end{remark}
